\section{Related Work}\label{sec:related}

Group Relative Policy Optimization (GRPO)~\cite{GRPO} addresses the computational overhead and training instability of PPO-style algorithms~\cite{PPO} by eliminating the need for a separate critic network. Instead, GRPO constructs baselines on-the-fly using multiple responses generated for each prompt. Specifically, GRPO samples a \emph{group} of multiple responses for each prompt and normalizes the rewards within this group to have zero mean and unit variance, creating relative advantages for policy updates.
However, this approach can be inefficient if all responses in a group receive the same reward (e.g., all incorrect or all correct), resulting in a zero-advantage for all samples and providing no learning signal.
To address this, DAPO~\cite{DAPO} enhances GRPO with engineering treatments like dynamic sampling, which continues generating responses until non-zero advantages are achieved, ensuring meaningful gradients.

Several other works have proposed improvements to group-based methods. \citet{zheng2025act} introduce GRESO, an online filtering algorithm that leverages reward training dynamics to predict and skip uninformative prompts \textit{before} generation. \citet{qu2025can} introduce a Bayesian estimation of the prompt accuracy and use it to form a bandit strategy, significantly reducing rollout overhead. \citet{liu2025part} propose ``Lite PPO'', which simplifies RLVR training to only advantage normalization and token-level loss aggregation.

Other group-based approaches include RLOO~\cite{RLOO}, which returns to the simpler REINFORCE~\cite{REINFORCE,rlbook} algorithm using a Leave-One-Out baseline that treats entire generations as single actions. Similarly, \citet{hao2025policy} propose On-Policy RL with Optimal Baseline (OPO), which uses a length-weighted average of rewards as an optimal simplified baseline. Despite these improvements, all group-based methods share fundamental limitations. They construct baselines from concurrently generated responses rather than persistent, historical estimates, inheriting the same core architectural constraints as GRPO: synchronization overhead and increased generation costs in distributed settings.

Moving beyond group-based methods, \citet{brantley2025accelerating} propose $A^*$-PO, a two-stage framework that achieves single-sample efficiency through a different approach. In the first stage, $A^*$-PO performs offline estimation to approximate the \textit{optimal} value function $V^*$ rather than the policy-specific value function $V_\pi$. The second stage uses this pre-computed optimal value to construct \emph{optimal} advantage estimates $A^*$ for a least-squares regression objective during online training. However, $A^*$-PO has key limitations compared to our approach. It relies on a \emph{fixed}, offline-computed estimate that does not adapt as the policy evolves during training. Additionally, $A^*$-PO is constrained by KL-regularized policy optimization, which restricts how far the optimized policy can deviate from the reference policy.

